% =============================================================================
% Chapter 7 -- Proposed Circular Oversampling Methods
% =============================================================================

\chapter{Proposed Circular Oversampling Methods}
\label{ch:proposed_methods}

This chapter presents the three proposed circular oversampling algorithms in full mathematical detail.
We begin with the common preprocessing pipeline shared by all methods, then present each algorithm with its formulation, pseudocode, and design rationale.


% ============================================================================
\section{Common Framework: Circle-Based Generation}
\label{sec:methods:framework}

All three algorithms share a common geometric framework based on four steps.

\subsection{Step 1: PCA Projection}

The minority class points $\{\vx_i\}_{i=1}^{n_{\min}} \subset \R^d$ are projected to a 2-D subspace via PCA:
\begin{equation}
    \hat{\vx}_i = (\vx_i - \bar{\vx}) \cdot \mathbf{V}_2, \quad \hat{\vx}_i \in \R^2,
    \label{eq:pca_proj}
\end{equation}
where $\bar{\vx}$ is the mean of the minority class and $\mathbf{V}_2 \in \R^{d \times 2}$ contains the first two principal component vectors.
This projection enables circular geometry in the 2-D subspace.
After generation, the inverse PCA transform (Equation~\eqref{eq:inverse_pca}) maps the 2-D disk back to $\R^d$ as a hyper-ellipsoidal region, since the transform scales coordinates by the corresponding PCA eigenvalues.
The circular shape is exact only in the projected subspace; in the original feature space the generation region is ellipsoidal, with aspect ratios determined by the ratio of PCA eigenvalues.

\subsection{Step 2: $k$-NN Graph Construction}

A $k$-nearest-neighbor graph is constructed on the projected minority points $\hat{\mathcal{X}}_{\min} = \{\hat{\vx}_i\}$.
For each point $\hat{\vx}_i$, the indices of its $k$ nearest neighbors in $\hat{\mathcal{X}}_{\min}$ are stored.

\subsection{Step 3: Circle Formation}

For each synthetic sample, a seed point $\hat{\vx}_i$ and a neighbor $\hat{\vx}_j$ are selected.
The circle is defined by their midpoint and half-distance:
\begin{align}
    \vc_{ij} &= \frac{\hat{\vx}_i + \hat{\vx}_j}{2}, \label{eq:circle_center}\\
    R_{ij}   &= \frac{\|\hat{\vx}_i - \hat{\vx}_j\|}{2}. \label{eq:circle_radius}
\end{align}
The resulting circle $\mathcal{D}(\vc_{ij}, R_{ij})$ is the closed disk centered at $\vc_{ij}$ with radius $R_{ij}$.

\subsection{Step 4: Inverse PCA Transform}

After generation in 2-D PCA space, synthetic points are lifted back to $\R^d$:
\begin{equation}
    \vx_{\text{new}} = \hat{\vx}_{\text{new}} \cdot \mathbf{V}_2^\top + \bar{\vx}.
    \label{eq:inverse_pca}
\end{equation}

\subsection{Native-Dimension Mode (\texttt{use\_pca=False})}
\label{sec:methods:nopca}

As an alternative to the 2-D PCA projection described above, all three algorithms support a \textbf{native-dimension mode} controlled by the \texttt{use\_pca=False} flag.
In this mode, Steps~1 and~4 are omitted entirely: all geometric operations (circle formation, ball sampling, clustering, and the $k$-NN graph) are performed directly in the original $d$-dimensional feature space.

\paragraph{Motivation.}
The PCA projection concentrates the geometry into the two directions of maximum variance, retaining between 52\% and 96\% of minority-class variance across the 14 benchmark datasets (median: 74\%).
The remaining variance — up to 48\% for the 34-feature ionosphere dataset — is discarded before any circle is formed, and synthetic points are subsequently lifted back via the inverse transform.

When \texttt{use\_pca=False}, no information is discarded.
The circle in Step~3 becomes a $d$-dimensional hyperball:
\begin{align}
    \vc_{ij} &= \tfrac{1}{2}(\vx_i + \vx_j), \\
    R_{ij}   &= \tfrac{1}{2}\|\vx_i - \vx_j\|_2.
\end{align}
The uniform sampling formula generalises to $d$-dimensional volume-uniform sampling:
\begin{equation}
    r = R_{ij} \cdot U^{1/d}, \quad U \sim \mathcal{U}(0, 1),
    \label{eq:nd_radius_sample}
\end{equation}
ensuring that the density of sampled points is uniform with respect to the $d$-dimensional volume element.
Gravity direction bias (GVM-CO, Algorithm~1) generalises from Von Mises angular sampling to the von Mises--Fisher (vMF) distribution on $\mathcal{S}^{d-1}$:
\begin{equation}
    \vz \sim \mathcal{N}(\bm{0}, \mathbf{I}_d), \quad \hat{\vz} = \frac{\vz + \kappa\,\hat{\bm{\mu}}}{\|\vz + \kappa\,\hat{\bm{\mu}}\|},
    \label{eq:vmf_approx}
\end{equation}
where $\hat{\bm{\mu}}$ is the unit-vector direction from the circle center toward the cluster gravity center and $\kappa$ is the same concentration score used in the 2-D mode.

\paragraph{Trade-off.}
The native-dimension mode avoids PCA information loss but introduces the \emph{curse of dimensionality}: in high-dimensional feature spaces ($d \gtrsim 15$), the volume of a $d$-dimensional ball concentrates near its surface, so uniform ball sampling produces points that are all approximately at radius $R$.
An empirical ablation comparing both modes is provided in Chapter~\ref{ch:ablation}.

The \textbf{Circular-SMOTE} baseline uses uniform disk sampling: $\theta \sim \mathcal{U}(0,2\pi)$, $r = R_{ij}\sqrt{U}$ ($U \sim \mathcal{U}(0,1)$), $\hat{\vx}_{\text{new}} = \vc_{ij} + r(\cos\theta,\sin\theta)^\top$.
The $\sqrt{U}$ radial scaling ensures area-uniform coverage (since the probability of a point at radius $r$ must be proportional to $r$).


% ============================================================================
\section{Gravity Scoring Framework}
\label{sec:methods:gravity}

Algorithms~1 and~3 use a gravity scoring framework that quantifies the attractiveness of each minority cluster.

\subsection{Cluster Metrics}

After clustering the projected minority points into $K$ groups $\{\mathcal{C}_1, \ldots, \mathcal{C}_K\}$, the following metrics are computed for each cluster $\mathcal{C}_c$ with $n_c$ points:

\begin{definition}[Cluster Spread]
\begin{equation}
    \mathrm{Spread}_c = \frac{1}{n_c} \sum_{i \in \mathcal{C}_c} \|\hat{\vx}_i - \boldsymbol{\mu}_c\|,
    \label{eq:spread}
\end{equation}
where $\boldsymbol{\mu}_c = \frac{1}{n_c}\sum_{i \in \mathcal{C}_c} \hat{\vx}_i$ is the cluster centroid.
\end{definition}

\begin{definition}[Cluster Area]
\begin{equation}
    \mathrm{Area}_c = \pi \cdot \mathrm{Spread}_c^2 + \epsilon.
    \label{eq:area}
\end{equation}
\end{definition}

\begin{definition}[Cluster Density]
\begin{equation}
    \mathrm{Density}_c = \frac{n_c}{\mathrm{Area}_c}.
    \label{eq:density}
\end{equation}
\end{definition}

\begin{definition}[Cluster Sparsity]
\begin{equation}
    \mathrm{Sparsity}_c = \frac{1}{n_c} \sum_{i \in \mathcal{C}_c} \frac{1}{k_{\text{nn}}} \sum_{j=1}^{k_{\text{nn}}} \|\hat{\vx}_i - \hat{\vx}_j^{(\text{nn})}\|,
    \label{eq:sparsity}
\end{equation}
where the inner sum is over the $k_{\text{nn}}$ nearest neighbors of $\hat{\vx}_i$ within cluster $\mathcal{C}_c$.
\end{definition}

\begin{definition}[Gravity Score]
\begin{equation}
    g_c = \frac{\mathrm{Density}_c}{\mathrm{Spread}_c \cdot \mathrm{Sparsity}_c + \epsilon}.
    \label{eq:gravity}
\end{equation}
The numerator $\mathrm{Density}_c$ measures how densely populated the cluster is per unit area.
The denominator $\mathrm{Spread}_c \cdot \mathrm{Sparsity}_c$ penalises clusters that are simultaneously spatially extended (large spread) \emph{and} internally loose (high mean nearest-neighbor distance).
A cluster that is compact and well-populated receives the highest gravity scores.
\end{definition}

To ensure scale-invariance, the raw gravity scores are normalised within each dataset:
\begin{equation}
    \tilde{g}_c = \frac{g_c - \min_{c'} g_{c'}}{\max_{c'} g_{c'} - \min_{c'} g_{c'} + \varepsilon},
    \label{eq:gravity_norm}
\end{equation}
so that $\tilde{g}_c \in [0,1]$ regardless of the absolute scale of the feature space.

\subsection{Gravity Center}

The gravity center is a density-weighted centroid pulled toward regions of high local density:

\begin{definition}[Gravity Center]
\begin{equation}
    \vg_c = \frac{\sum_{i \in \mathcal{C}_c} w_i \cdot \hat{\vx}_i}{\sum_{i \in \mathcal{C}_c} w_i}, \quad w_i = \frac{1}{\sigma_i^{\text{local}} + \epsilon},
    \label{eq:gravity_centre}
\end{equation}
where $\sigma_i^{\text{local}}$ is the mean $k$-NN distance of point $\hat{\vx}_i$ within its cluster.
Points in denser local regions have smaller $\sigma_i^{\text{local}}$ and hence larger weights $w_i$, pulling the gravity center toward the densest part of the cluster.
\end{definition}


% ============================================================================
\section{Algorithm 1: Gravity-Biased Von Mises Circular Oversampling (GVM-CO)}
\label{sec:methods:gvm_co}

GVM-CO generates synthetic samples inside circles formed between minority point pairs, with the angular distribution biased by a Von Mises distribution directed toward the gravity center of the seed's cluster.

\subsection{Von Mises Direction}

For a circle defined by seed $\hat{\vx}_i$ (in cluster $c$) and neighbor $\hat{\vx}_j$, with center $\vc_{ij}$, the Von Mises mean direction is the angle from the circle center toward the gravity center:
\begin{equation}
    \mu = \arctan2\!\left(\vg_c^{(y)} - \vc_{ij}^{(y)},\; \vg_c^{(x)} - \vc_{ij}^{(x)}\right).
    \label{eq:mu_direction}
\end{equation}

\subsection{Von Mises Concentration}

The concentration $\kappa$ controls the strength of directional biasing and is computed from two sub-scores.

\paragraph{Distance score $s_d$.} Normalised distance from the circle center to the gravity center (closer centers get higher scores):
\begin{equation}
    s_d = \frac{d_{\max} - d(\vc_{ij}, \vg_c)}{d_{\max} - d_{\min} + \epsilon},
    \label{eq:s_d}
\end{equation}
where $d_{\min}$ and $d_{\max}$ are the extreme distances across all circles in the current batch.

\paragraph{Gravity score $s_g$.} Normalised gravity score of the seed's cluster:
\begin{equation}
    s_g = \frac{g_c - g_{\min}}{g_{\max} - g_{\min} + \epsilon}.
    \label{eq:s_g}
\end{equation}

\paragraph{Combined score.} The two sub-scores are blended using an additive--multiplicative combination:
\begin{equation}
    s_{\text{combined}} = \alpha \cdot \frac{s_d + s_g}{2} + (1 - \alpha) \cdot s_d \cdot s_g,
    \label{eq:s_combined}
\end{equation}
where $\alpha \in [0, 1]$ controls the balance between additive (robust) and multiplicative (selective) combination.

\paragraph{Concentration mapping.}
\begin{equation}
    \kappa = \kappa_{\max} \cdot s_{\text{combined}}^{\gamma},
    \label{eq:kappa}
\end{equation}
where $\kappa_{\max}$ is the maximum concentration and $\gamma > 0$ is a nonlinearity exponent.
A minimum floor prevents complete deactivation of the bias:
\begin{equation}
    \kappa_c = \max\!\left(\kappa_{\max} \cdot s_c^{\gamma},\; \varepsilon\right),
    \label{eq:kappa_floor}
\end{equation}
where $\varepsilon = 10^{-6}$.

\subsection{Sampling Procedure}

Each synthetic point is generated as:
\begin{align}
    \theta &\sim \mathrm{VonMises}(\mu, \kappa), \label{eq:theta_gvm}\\
    r      &= R_{ij} \cdot \sqrt{U}, \quad U \sim \mathcal{U}(0,1), \label{eq:r_gvm}\\
    \hat{\vx}_{\text{new}} &= \vc_{ij} + r \cdot (\cos\theta,\; \sin\theta)^{\top}. \label{eq:x_new_gvm}
\end{align}

\subsection{Cross-Cluster Border Restriction}

An optional variant of GVM-CO (GVM-CO-CC) enforces that the neighbor $\hat{\vx}_j$ comes from the nearest \emph{other} cluster:
\begin{equation}
    c^* = \argmin_{c' \neq c} \|\boldsymbol{\mu}_c - \boldsymbol{\mu}_{c'}\|.
    \label{eq:nearest_cluster}
\end{equation}
The $k$-NN search for neighbors of seeds in cluster $c$ is restricted to points in cluster $c^*$.
This promotes generation in inter-cluster border regions, improving connectivity between minority sub-populations.
As shown in the experimental results (Chapter~\ref{ch:results}), GVM-CO-CC is most useful at low imbalance ratios; at high imbalance the inter-cluster region is typically majority-dominated, making cross-cluster pairing counterproductive.

\subsection{Full-Dataset Clustering Variant}

Inspired by K-Means SMOTE~\cite{douzas2018improving}, a third variant applies K-means to the \emph{full dataset} $\mathcal{X} = \mathcal{X}_{\min} \cup \mathcal{X}_{\text{maj}}$.
The procedure:
\begin{enumerate}[leftmargin=*]
    \item Apply K-means with $K$ clusters to the full projected dataset.
    \item For each cluster $c$, compute the minority ratio:
    \begin{equation}
        \rho_c = \frac{|\{i : c_i = c, \; y_i = y_{\min}\}|}{|\{i : c_i = c\}|}.
    \end{equation}
    \item Select clusters where $\rho_c \geq \rho_{\min}$ (default $\rho_{\min} = \text{IR}^{-1}$).
    \item Within each selected cluster, compute gravity scores and centers using only the minority points.
    \item Apply GVM-CO generation within selected clusters, with synthetic samples allocated proportionally to the number of minority points per cluster.
\end{enumerate}
The key advantage is that cluster boundaries are informed by the majority distribution, naturally avoiding generation in majority-dominated regions.

\begin{algorithm}[H]
\caption{GVM-CO: Gravity-biased Von Mises Circular Oversampling}
\label{alg:gvm_co}
\KwInput{Minority $\mathcal{X}_{\min}$, $K$, $k_{\text{seed}}$, $\kappa_{\max}$, $\gamma$, $\alpha$}
\KwOutput{Synthetic points $\mathcal{S}$}

$\hat{\mathcal{X}}_{\min} \leftarrow \mathrm{PCA}(\mathcal{X}_{\min}, p=2)$\;
$\{c_i\} \leftarrow \mathrm{K\text{-}Means}(\hat{\mathcal{X}}_{\min}, K)$\;
Compute gravity scores $\{g_c\}$ and gravity centers $\{\vg_c\}$ via Eqs.~\eqref{eq:gravity}--\eqref{eq:gravity_centre}\;
Build $k_{\text{seed}}$-NN graph on $\hat{\mathcal{X}}_{\min}$\;
$\mathcal{S} \leftarrow \emptyset$\;
\For{$t = 1$ \KwTo $n_{\text{synth}}$}{
    Select random seed $\hat{\vx}_i$ in cluster $c$\;
    Select random neighbor $\hat{\vx}_j$ from $k$-NN of $\hat{\vx}_i$\;
    $\vc \leftarrow \frac{1}{2}(\hat{\vx}_i + \hat{\vx}_j)$;\quad $R \leftarrow \frac{1}{2}\|\hat{\vx}_i - \hat{\vx}_j\|$\;
    $\mu \leftarrow \arctan2(\vg_c - \vc)$\;
    Compute $s_d$, $s_g$, $s_{\text{combined}}$, $\kappa$ via Eqs.~\eqref{eq:s_d}--\eqref{eq:kappa_floor}\;
    $\theta \sim \mathrm{VonMises}(\mu, \kappa)$;\quad $r \leftarrow R\sqrt{U}$, $U \sim \mathcal{U}(0,1)$\;
    $\hat{\vx}_{\text{new}} \leftarrow \vc + r(\cos\theta, \sin\theta)^\top$\;
    $\mathcal{S} \leftarrow \mathcal{S} \cup \{\hat{\vx}_{\text{new}}\}$\;
}
Lift $\mathcal{S}$ to $\R^d$ via inverse PCA (Eq.~\eqref{eq:inverse_pca})\;
\Return{$\mathcal{S}$}
\end{algorithm}


% ============================================================================
\section{Algorithm 2: Local Region Estimation Circular Oversampling (LRE-CO)}
\label{sec:methods:lre_co}

LRE-CO addresses cross-boundary contamination by restricting generation to the intersection of the circular disk with a Voronoi cell defined by the cluster partition.

\subsection{Voronoi Partitioning}

Given $K$ cluster centroids $\{\boldsymbol{\mu}_c\}_{c=1}^K$, the 2-D space is partitioned into Voronoi cells:
\begin{equation}
    \mathcal{V}_c = \left\{\hat{\vx} \in \R^2 \;\middle|\; \|\hat{\vx} - \boldsymbol{\mu}_c\| \leq \|\hat{\vx} - \boldsymbol{\mu}_{c'}\|,\; \forall c' \neq c\right\}.
    \label{eq:voronoi}
\end{equation}
Each point in $\R^2$ belongs to exactly one Voronoi cell.

\subsection{Certainty Threshold}

Before generating synthetic samples within a circle, LRE-CO evaluates whether the circle lies in a minority-dominated region.
The certainty of a disk $\mathcal{D}(\vc_{ij}, R_{ij})$ is defined as:
\begin{equation}
    \mathrm{certainty}(\mathcal{D}) = \frac{n_{\min}^{\text{inside}}}{n_{\min}^{\text{inside}} + n_{\text{maj}}^{\text{inside}}},
    \label{eq:certainty}
\end{equation}
where $n_{\min}^{\text{inside}}$ and $n_{\text{maj}}^{\text{inside}}$ are the number of minority and majority points inside the disk, respectively.
If $\mathrm{certainty}(\mathcal{D}) < \tau$ (default $\tau = 0.80$), the circle is rejected, preventing synthetic generation in majority-dominated regions.

\begin{remark}[2-D Approximation]
The certainty metric is computed in the 2-D PCA subspace.
For datasets where the first two principal components explain a small fraction of the total variance, a disk that appears minority-dominated in 2-D may straddle a class boundary in the remaining dimensions.
The effectiveness of the certainty mechanism therefore diminishes as the intrinsic dimensionality of the minority distribution increases relative to the PCA subspace.
Practitioners should verify the explained variance of the 2-D projection and consider using a stricter threshold $\tau$ for high-dimensional datasets (see Table~\ref{tab:setup:pca_variance} in Chapter~\ref{ch:experimental_setup}).
\end{remark}

\subsection{Constrained Generation}

For an eligible circle defined by seed $\hat{\vx}_i$ (in cluster $c$) and neighbor $\hat{\vx}_j$, LRE-CO restricts synthetic point generation to:
\begin{equation}
    \hat{\vx}_{\text{new}} \in \mathcal{D}(\vc_{ij}, R_{ij}) \cap \mathcal{V}_c.
    \label{eq:lre_constraint}
\end{equation}

\subsection{Rejection Sampling}

Since the intersection $\mathcal{D}(\vc_{ij}, R_{ij}) \cap \mathcal{V}_c$ does not have a simple closed-form shape, LRE-CO uses rejection sampling:
\begin{enumerate}[leftmargin=*]
    \item Generate a candidate point $\hat{\vx}_{\text{cand}}$ uniformly inside the disk $\mathcal{D}(\vc_{ij}, R_{ij})$.
    \item Check whether $\hat{\vx}_{\text{cand}} \in \mathcal{V}_c$ by verifying $\argmin_{c'} \|\hat{\vx}_{\text{cand}} - \boldsymbol{\mu}_{c'}\| = c$.
    \item If accepted, add to the synthetic set. Otherwise reject and repeat.
\end{enumerate}
The acceptance rate depends on the overlap between the disk and the Voronoi cell.
For circles whose center is near the cluster centroid, the acceptance rate is typically high; a maximum iteration limit prevents infinite loops.

\subsection{Allocation Strategy}

Synthetic samples are allocated to clusters proportionally to cluster size weighted by normalised gravity score:
\begin{equation}
    n_c^{\text{synth}} = n_{\text{synth}} \cdot \frac{n_c \cdot \tilde{g}_c}{\sum_{c'=1}^{K} n_{c'} \cdot \tilde{g}_{c'}}.
    \label{eq:lre_allocation}
\end{equation}
This assigns more synthetic samples to clusters that are both large and high-gravity.
Users requiring uniform coverage across all minority sub-populations can set all $\tilde{g}_c = 1$, reducing the allocation to size-proportional.

\begin{algorithm}[H]
\caption{LRE-CO: Local Region Estimation Circular Oversampling}
\label{alg:lre_co}
\KwInput{Minority $\mathcal{X}_{\min}$, majority $\mathcal{X}_{\text{maj}}$, $K$, $k_{\text{seed}}$, $\kappa_{\max}$, $\gamma$, $\alpha$, $\tau$}
\KwOutput{Synthetic points $\mathcal{S}$}

$\hat{\mathcal{X}}_{\min} \leftarrow \mathrm{PCA}(\mathcal{X}_{\min}, p=2)$\;
$\{c_i\}, \{\boldsymbol{\mu}_c\} \leftarrow \mathrm{K\text{-}Means}(\hat{\mathcal{X}}_{\min}, K)$\;
Compute gravity scores, gravity centers\;
Compute Voronoi partition $\{\mathcal{V}_c\}_{c=1}^K$ via Eq.~\eqref{eq:voronoi}\;
Allocate $n_c^{\text{synth}}$ per cluster via Eq.~\eqref{eq:lre_allocation}\;
$\mathcal{S} \leftarrow \emptyset$\;
\For{each cluster $c = 1, \ldots, K$}{
    \For{$t = 1$ \KwTo $n_c^{\text{synth}}$}{
        Select seed $\hat{\vx}_i \in \mathcal{C}_c$, neighbor $\hat{\vx}_j$\;
        Form circle $(\vc, R)$ via Eqs.~\eqref{eq:circle_center}--\eqref{eq:circle_radius}\;
        \If{$\mathrm{certainty}(\mathcal{D}) < \tau$ (Eq.~\eqref{eq:certainty})}{
            Skip circle and select new pair\;
        }
        \Repeat{$\hat{\vx}_{\text{cand}} \in \mathcal{V}_c$ \textbf{or} max iterations}{
            $\theta \sim \mathcal{U}(0, 2\pi)$;\quad $r \leftarrow R\sqrt{U}$\;
            $\hat{\vx}_{\text{cand}} \leftarrow \vc + r(\cos\theta, \sin\theta)^\top$\;
        }
        $\mathcal{S} \leftarrow \mathcal{S} \cup \{\hat{\vx}_{\text{cand}}\}$\;
    }
}
Lift $\mathcal{S}$ to $\R^d$ via inverse PCA\;
\Return{$\mathcal{S}$}
\end{algorithm}


% ============================================================================
\section{Algorithm 3: Layered Segmental Circular Oversampling (LS-CO)}
\label{sec:methods:ls_co}

LS-CO introduces a structured decomposition of the circular region into concentric annular layers and angular segments, with differentiated sampling densities.

\subsection{Layer Decomposition}

Each circle of radius $R_{ij}$ is divided into $L$ concentric annular layers with equal radial width $\Delta r = R_{ij}/L$:
\begin{equation}
    \mathcal{A}_\ell = \left\{\hat{\vx} \;\middle|\; r_{\ell-1} \leq \|\hat{\vx} - \vc_{ij}\| \leq r_\ell\right\}, \quad r_\ell = \frac{\ell}{L} \cdot R_{ij}, \quad \ell = 1, \ldots, L.
    \label{eq:layers}
\end{equation}
The rings have equal radial width but unequal areas: the area of ring $\ell$ is $\pi(r_\ell^2 - r_{\ell-1}^2) = \pi R_{ij}^2 (2\ell-1)/L^2$, so the outermost ring is $(2L-1)$ times larger than the innermost.
Synthetic samples are allocated uniformly across layers by default, deliberately over-sampling the inner rings relative to their area, concentrating generation near the circle center where the seed minority instances are located and the minority class is most reliably represented.

\subsection{Angular Segmentation}

Each layer is divided into $M$ angular segments:
\begin{equation}
    \mathcal{S}_{\ell,m} = \mathcal{A}_\ell \cap \left\{\hat{\vx} \;\middle|\; \theta_{m-1} \leq \angle(\hat{\vx} - \vc_{ij}) < \theta_m\right\}, \quad \theta_m = \frac{2\pi m}{M}.
    \label{eq:segments}
\end{equation}

\subsection{Layer Density Allocation}

The number of synthetic samples in each layer is determined by a gravity-proximity weight:
\begin{equation}
    w_\ell = \exp\!\left(-\beta \cdot \frac{|r_\ell^* - r_{\text{gc}}|}{R_{ij}}\right),
    \label{eq:layer_weight}
\end{equation}
where $r_\ell^* = \frac{r_{\ell-1} + r_\ell}{2}$ is the midpoint radius and $r_{\text{gc}}$ is the distance from the circle center to the projection of the gravity center onto the radial axis.
The allocation per layer is:
\begin{equation}
    n_\ell = n_{\text{per-circle}} \cdot \frac{w_\ell}{\sum_{\ell'} w_{\ell'}}.
    \label{eq:layer_alloc}
\end{equation}

\subsection{Within-Segment Sampling}

Within each segment $\mathcal{S}_{\ell,m}$, the angular distribution uses a Von Mises distribution with segment-local concentration:
\begin{equation}
    \kappa_{\ell,m} = \kappa_{\max} \cdot \frac{L - \ell + 1}{L} \cdot \exp\!\left(-\frac{|\theta_m^{\text{mid}} - \mu|^2}{2\sigma_\theta^2}\right),
    \label{eq:seg_kappa}
\end{equation}
where $\theta_m^{\text{mid}} = \frac{\theta_{m-1} + \theta_m}{2}$ is the midpoint angle and $\mu$ is the gravity direction.
Inner layers (smaller $\ell$) and segments aligned with $\mu$ receive higher concentration.

The radial component within layer $\ell$ uses area-corrected uniform sampling:
\begin{equation}
    r = \sqrt{r_{\ell-1}^2 + U \cdot (r_\ell^2 - r_{\ell-1}^2)}, \quad U \sim \mathcal{U}(0, 1).
    \label{eq:annulus_sample}
\end{equation}

\begin{algorithm}[H]
\caption{LS-CO: Layered Segmental Circular Oversampling}
\label{alg:ls_co}
\KwInput{Minority $\mathcal{X}_{\min}$, $K$, $k_{\text{seed}}$, $L$, $M$, $\kappa_{\max}$, $\beta$, $\sigma_\theta$}
\KwOutput{Synthetic points $\mathcal{S}$}

$\hat{\mathcal{X}}_{\min} \leftarrow \mathrm{PCA}(\mathcal{X}_{\min}, p=2)$\;
$\{c_i\} \leftarrow \mathrm{Cluster}(\hat{\mathcal{X}}_{\min}, K)$\;
Compute gravity scores and centers via Eqs.~\eqref{eq:gravity}--\eqref{eq:gravity_centre}\;
Build $k_{\text{seed}}$-NN graph\;
$\mathcal{S} \leftarrow \emptyset$\;
\For{each synthetic sample}{
    Select seed $\hat{\vx}_i$ and neighbor $\hat{\vx}_j$\;
    Form circle $(\vc, R)$ via Eqs.~\eqref{eq:circle_center}--\eqref{eq:circle_radius}\;
    Compute gravity direction $\mu = \arctan2(\vg_c - \vc)$\;
    Compute layer weights $\{w_\ell\}_{\ell=1}^L$ (Eq.~\eqref{eq:layer_weight}) and allocations $\{n_\ell\}$ (Eq.~\eqref{eq:layer_alloc})\;
    \For{$\ell = 1$ \KwTo $L$}{
        \For{each sample allocated to layer $\ell$}{
            Select segment $m$ with Von Mises-weighted probability\;
            $\kappa_{\ell,m} \leftarrow$ Eq.~\eqref{eq:seg_kappa}\;
            $\theta \sim \mathrm{VonMises}(\mu, \kappa_{\ell,m})$ restricted to $[\theta_{m-1}, \theta_m)$\;
            $r \leftarrow$ area-corrected sample in $[r_{\ell-1}, r_\ell]$ (Eq.~\eqref{eq:annulus_sample})\;
            $\hat{\vx}_{\text{new}} \leftarrow \vc + r(\cos\theta, \sin\theta)^\top$\;
            $\mathcal{S} \leftarrow \mathcal{S} \cup \{\hat{\vx}_{\text{new}}\}$\;
        }
    }
}
Lift $\mathcal{S}$ to $\R^d$ via inverse PCA\;
\Return{$\mathcal{S}$}
\end{algorithm}

LS-CO has two primary variants: \textbf{LS-CO (G)} uses a single gravity center computed from all minority points, and \textbf{LS-CO (C)} uses per-cluster gravity centers computed after K-means clustering, providing finer control for multi-modal minority distributions.


% ============================================================================
\section{Optional Post-Hoc Denoising}
\label{sec:methods:denoising}

All three algorithms support optional denoising after synthetic sample generation:

\begin{itemize}[leftmargin=*]
    \item \textbf{Tomek Links}: Identifies pairs of instances from different classes that are mutual nearest neighbors and removes the majority instance (or both), cleaning the decision boundary.
    \item \textbf{Edited Nearest Neighbours (ENN)}: Removes any instance whose class label disagrees with the majority of its $k$-nearest neighbors. More aggressive than Tomek Links.
\end{itemize}

Denoising is applied to the \emph{combined} dataset (original + synthetic) before classifier training.
As shown in the ablation (Chapter~\ref{ch:ablation}), Tomek links provide a small net improvement for GVM-CO and LRE-CO, while ENN is counterproductive for these methods because the circular generation already incorporates boundary-avoidance mechanisms internally.


% ============================================================================
\section{Comparison of the Three Algorithms}
\label{sec:methods:comparison}

Table~\ref{tab:method_comparison} summarises the key differences between the three algorithms.

\begin{table}[H]
\centering
\caption{Comparison of the three proposed circular oversampling algorithms.}
\label{tab:method_comparison}
\small
\begin{tabular}{p{3.5cm}p{3.5cm}p{3.5cm}p{3.5cm}}
\toprule
\textbf{Feature} & \textbf{GVM-CO} & \textbf{LRE-CO} & \textbf{LS-CO} \\
\midrule
Angular distribution & Von Mises & Uniform (or Von Mises) + rejection & Von Mises per segment \\
Region constraint & None (full disk) & Voronoi cell & Layer + segment \\
Spatial control & Global (gravity direction) & Local (Voronoi boundary) & Structured (layers) \\
Cross-cluster variant & Yes (GVM-CO-CC) & N/A & N/A \\
Full-dataset variant & Yes (KM-SMOTE style) & N/A & N/A \\
Rejection sampling & No & Yes & No \\
Hyperparameters & $K, k, \kappa_{\max}, \gamma, \alpha$ & $K, k, \kappa_{\max}, \gamma, \alpha, \tau$ & $K, k, L, M, \kappa_{\max}, \beta, \sigma_\theta$ \\
Complexity (generation) & $\mathcal{O}(n_{\text{synth}})$ & $\mathcal{O}(n_{\text{synth}} \cdot K/p_{\text{accept}})$ & $\mathcal{O}(n_{\text{synth}} \cdot L)$ \\
Preprocessing & \multicolumn{3}{c}{$\mathcal{O}(n_{\min} \log n_{\min})$ (shared: PCA + $k$-NN graph)} \\
\bottomrule
\end{tabular}
\end{table}


% ============================================================================
\section{Computational Complexity}
\label{sec:methods:complexity}

The computational cost of each algorithm has five components:

\begin{enumerate}[leftmargin=*]
    \item \textbf{PCA}: $\mathcal{O}(n_{\min} \cdot d^2)$ for the covariance matrix and eigendecomposition. Negligible for moderate $d$ since $n_{\min}$ is small.
    \item \textbf{$k$-NN graph}: $\mathcal{O}(n_{\min}^2)$ naively, or $\mathcal{O}(n_{\min} \log n_{\min})$ with a KD-tree in the 2-D PCA space.
    \item \textbf{Clustering}: $\mathcal{O}(n_{\min} \cdot K \cdot T_{\text{iter}})$ for K-Means with $T_{\text{iter}}$ iterations.
    \item \textbf{Generation}: $\mathcal{O}(n_{\text{synth}})$ per sample for GVM-CO and LS-CO; $\mathcal{O}(n_{\text{synth}} / p_{\text{accept}})$ for LRE-CO, where $p_{\text{accept}}$ is the average fraction of disk area inside the Voronoi cell.
\end{enumerate}

Since $n_{\min}$ is typically small by definition, the dominant cost is usually the generation step.
All three methods are comparable in computational cost to K-Means SMOTE and are substantially faster than GAN-based approaches.


% ============================================================================
\section{Chapter Summary}
\label{sec:methods:summary}

This chapter has presented the three proposed circular oversampling algorithms.
All three share a common framework based on PCA projection, $k$-NN graph construction, circle formation, and inverse PCA lifting.
They differ in how they control the placement of synthetic samples within each circle:

\begin{enumerate}[leftmargin=*]
    \item \textbf{GVM-CO} combines clustering-based gravity scoring with Von Mises directional sampling, biasing generation toward dense minority regions while preserving diversity through the radial component.
    \item \textbf{LRE-CO} restricts generation to the intersection of circular disks with Voronoi cells, preventing cross-boundary contamination and enabling fine-grained local generation.
    \item \textbf{LS-CO} decomposes circular regions into layers and segments, with gravity-weighted density allocation and per-segment Von Mises concentration, providing structured control over sample placement.
\end{enumerate}

The next chapter describes the experimental setup used to evaluate these algorithms.
